\documentclass[../main.tex]{subfiles}

\begin{document}

\section{Interpolacao Polinomial}

\subsection{Polinomios de Lagrange}

A interpolacao de Lagrange foi implementada montando cada base \(L_i(x)\)
diretamente:
\[
    L_i(x) = \prod_{j \ne i} \frac{x - x_j}{x_i - x_j}.
\]
O polinomio final e calculado por
\[
    P(x) = \sum_{i=0}^{n} y_i L_i(x).
\]

No codigo, os termos sao multiplicados como polinomios simples, guardados pela
lista de coeficientes \([c_0,c_1,\ldots,c_n]\). Isso deixa facil avaliar o
polinomio em qualquer ponto e tambem escrever a expressao obtida no CSV.

\subsection{Diferencas divididas de Newton}

No metodo de Newton, primeiro e montada a tabela de diferencas divididas. O
polinomio fica na forma:
\[
    P(x) = f[x_0] + f[x_0,x_1](x-x_0) + \cdots
    + f[x_0,\ldots,x_n](x-x_0)\cdots(x-x_{n-1}).
\]

Essa forma e conveniente porque permite avaliar o polinomio com uma estrutura
parecida com o metodo de Horner. Para facilitar a leitura dos resultados, o
codigo tambem converte o polinomio para potencias de \(x\).

\subsection{Problemas teste}

\subsubsection{Exercicio 10.2: trajetoria de um projetil}

O exercicio 10.2 fornece os pontos \((0,0)\), \((10,6)\) e \((20,4)\) da
trajetoria de um projetil. Usando interpolacao por Lagrange, foi obtido:
\[
    P(x) = x - 0.04x^2.
\]

Comparando com a equacao teorica
\[
    y = x\tan\psi - \frac{1}{2}g\frac{x^2}{v_0^2\cos^2\psi},
\]
temos \(\tan\psi = 1\), logo \(\psi = 45^\circ\). Com \(g=9.86\), a velocidade
inicial calculada foi:
\[
    v_0 = 15.7003 \text{ m/s}.
\]
No ponto \(x=5\), a altitude calculada foi:
\[
    P(5)=4.0\text{ m}.
\]
Os resultados estao em \path{output/exercicios/ex10_2.csv}.

\subsubsection{Exercicio 10.6: resistencia de um fio}

O exercicio 10.6 pede estimativas de resistencia para comprimentos de 1730 m e
3200 m usando parabolas de grau 2 e 3. Para cada consulta, o programa escolheu
os pontos mais proximos e montou o polinomio interpolador por diferencas
divididas de Newton.

Os valores obtidos foram:
\[
\begin{array}{c|c|c}
\text{comprimento} & \text{grau} & \text{resistencia aproximada} \\
\hline
1730 & 2 & 9.241544 \\
1730 & 3 & 9.292921 \\
3200 & 2 & 17.796800 \\
3200 & 3 & 17.899840
\end{array}
\]
A pequena diferenca entre grau 2 e grau 3 mostra o efeito de incluir mais um
ponto na interpolacao. A saida completa esta em
\path{output/exercicios/ex10_6.csv}.

\subsubsection{Exercicio 10.9: circuito com varistor}

No exercicio 10.9, a corrente \(I\) e tabelada em funcao da tensao \(E_2\) no
varistor. No item (a), foi usado o polinomio interpolador sobre todos os pontos
para calcular \(I\) quando \(E_2=2.3\). O valor obtido foi:
\[
    I(2.3) = 0.810152.
\]
Como \(R_1=10\Omega\), entao:
\[
    E_1 = R_1 I = 8.10152,
    \quad
    E = E_1 + E_2 = 10.40152.
\]

No item (b), foi usada interpolacao inversa de grau 2. Para \(E_1=10\), temos
\(I=1\), e a tensao estimada no varistor foi:
\[
    E_2 = 2.46417\text{ V}.
\]
Os resultados estao registrados em \path{output/exercicios/ex10_9.csv}.

\end{document}
